\documentclass[10pt,a4paper]{amsart}
\usepackage[margin=19mm]{geometry}
\usepackage{amsmath,amssymb,mathtools}
\usepackage[scr=rsfs,cal=euler]{mathalfa}
\usepackage{xfrac}
\usepackage[unicode,hidelinks,bookmarks=false]{hyperref}
\newcommand{\ie}{i.e.\ }
\newcommand{\cf}{cf.\ }
\newcommand{\resp}{respectively\ }
\newcommand{\Subsection}{\subsection}
\providecommand{\nobreakdash}{\nobreak\hbox{-}}
\setlength{\emergencystretch}{2em}
\multlinegap0pt
\usepackage{bm}
\usepackage{bbm}
\usepackage{dsfont}
\usepackage{xspace}
\usepackage{cancel}
\usepackage{mathtools}
\usepackage{amsthm}
\makeatletter
\@ifundefined{theorem}{\newtheorem{theorem}{Theorem}}{}
\@ifundefined{lemma}{\newtheorem{lemma}[theorem]{Lemma}}{}
\@ifundefined{proposition}{\newtheorem{proposition}[theorem]{Proposition}}{}
\makeatother
\newcommand{\EE}{\mathbb{E}}
\newcommand{\PP}{\mathbb{P}}
\title[The Frog Model on $\mathbb{Z}$ with Discrete Weibull Lifetim]{The Frog Model on $\mathbb{Z}$ with Discrete Weibull Lifetimes and Random Parameter $p$}
\author{J. H. Ram\'irez Gonz\'alez, Gustavo O. Carvalho, F\'abio P. Machado}
\date{Source: arXiv:2601.15526v1 (CC BY 4.0)}
\begin{document}
\maketitle

\noindent\emph{Source and licence.} The Frog Model on $\mathbb{Z}$ with Discrete Weibull Lifetimes and Random Parameter $p$. Version arXiv:2601.15526v1 (CC BY 4.0), distributed under the Creative Commons Attribution 4.0 International licence, as recorded in the arXiv licence metadata for this exact version and shown on the abstract page https://arxiv.org/abs/2601.15526v1. Full rights notice: licenses/arxiv-2601.15526-CC-BY-4.0.txt.\par\smallskip

\noindent\textit{Editorial scope.} Counted unit: the Lemma whose source label is \texttt{lem:tauber-Tgamma-EN-eps} in the article (source0.tex lines 1096--1105, source label \texttt{lem:tauber-Tgamma-EN-eps}), delivered together with the article's own complete original proof of it (source0.tex lines 568--639, one \texttt{proof} environment of 72 lines). No mathematics is written, repaired or re-derived: statement and proof are the article's own text line for line. The proof is detached from the statement in the source: the article states the result at lines 1096--1105 and prints its proof at lines 568--639, and the proof environment's own header names the statement, so the association is the article's own. Required context retained uncounted: eq:edge-behavior (lines 165--167); prop:upper-gamma>1-EN-eps (lines 469--484); eq:upper-gamma>1-EN-eps (lines 471--473); sumindp (lines 542--546); eq:edge-band (lines 551--555); lem:potter-EN-eps (lines 1031--1040); eq:potter-EN-eps (lines 1033--1035); eq:UCT-EN-eps (lines 1037--1039); lem:convex-ineq-EN-eps (lines 1134--1145). Every definition, display and label of those lines is retained unchanged. Omitted from this package and not redistributed: 1--164, 168--468, 474--541, 547--550, 556--567, 640--1030, 1036--1036, 1040--1095, 1106--1133, 1146--1364. Nothing inside a retained excerpt is omitted or altered: every retained body is delivered exactly as the article's lines are written, so no omission receipt of any kind accompanies this package. Citations: the retained excerpts carry no citation command at all, so no bibliography entry of the article is transcribed and no reference is redistributed. Reference adaptation: none was needed and none was made; every reference inside the delivered unit resolves inside it, so no unresolved reference or citation is produced. Printed slips kept literal: no printed word, symbol, hypothesis or number of the counted statement or of its proof was altered. Layout and class: the article's own class is \texttt{imsart} with packages including amsthm, amsmath, amsfonts, amssymb, natbib, hyperref, graphicx, xcolor; the wrapper is the lane's pinned \texttt{amsart} editorial wrapper at 10pt on A4, which restates the theorem-like environments the retained text needs and adds the article's own macro definitions that the retained text needs (\texttt{\textbackslash EE}, \texttt{\textbackslash PP}), with extra wrapper packages (bm, bbm, dsfont, xspace, cancel, mathtools). The delivered numbering is the unit's own. Assets: no retained span contains a figure environment, a graphics-inclusion command, a PDF-page inclusion, a table or a TikZ picture; no image file is copied into this package, the package declares no asset dependency and no image file is redistributed with it. Licence: version arXiv:2601.15526v1 of this article is distributed under the Creative Commons Attribution 4.0 International licence, as recorded in the arXiv licence metadata for this exact version and shown on the abstract page https://arxiv.org/abs/2601.15526v1. A full rights notice is delivered at licenses/arxiv-2601.15526-CC-BY-4.0.txt. Characters: the retained excerpts contain 1 characters outside ASCII, all of them Latin-1 letters that the delivered engine XeLaTeX renders with its default Latin Modern text font with no missing character.
\begin{equation}\label{eq:edge-behavior}
f_\pi(u)\ \sim\ (1-u)^{\beta-1}\,L\!\Big(\frac{1}{1-u}\Big)\qquad (u\uparrow 1),
\end{equation}

\begin{proposition}[Upper bound \(\gamma > 1\)]\label{prop:upper-gamma>1-EN-eps}
Assume \(\gamma>1\) and \eqref{eq:edge-behavior}. Then, for every $\varepsilon\in(0,(2\gamma\beta)\wedge\tfrac{1}{2})$ and \(C>1\) there exists $N_{\varepsilon,C}$ such that for all $n\ge N_{\varepsilon,C}$,
\begin{equation}\label{eq:upper-gamma>1-EN-eps}
n\,\PP(D^{\rightarrow}\ge n)\ \le\ n^{1-2\gamma\beta}L(n^{2\gamma})D^+_{\gamma,\beta,\varepsilon,C}+n2^{-n^\gamma},
\end{equation}


where 
\[
D^+_{\gamma,\beta,\epsilon,C}:=2\gamma (1+\varepsilon)^{2} c_{\varepsilon}^{-2\gamma\beta}C\int_{0}^{\infty}e^{-y}y^{2\gamma\beta-1}\max\{y^{\varepsilon},y^{-\varepsilon}\}dy
\]

and \(c_{\varepsilon}:=(1-\varepsilon)\Gamma\!\bigl(1-\tfrac{1}{2\gamma}\bigr)\sqrt{\tfrac{2}{\pi}}2^{-1/(2\gamma)}\). 


\end{proposition}

\begin{equation}\label{sumindp}
\EE[s^{\tau_n}]=f(s)^n:=\Big(\tfrac{1-\sqrt{1-s^2}}{s}\Big)^{\!n}
\quad\text{and}\quad
\tau_n\ \stackrel{d}{=}\ \sum_{i=1}^{n} \tau_1^{i},\ \ \ \tau_1^{i}\ \text{i.i.d. with }\tau_1^{i}\stackrel{d}{=}\tau_1,
\end{equation}

\begin{equation}\label{eq:edge-band}
(1-\varepsilon)\,h^{\beta-1}L(1/h)
\le f_\pi(1-h)\le
(1+\varepsilon)\,h^{\beta-1}L(1/h).
\end{equation}

\begin{lemma}\label{lem:potter-EN-eps}
If $L$ is slowly varying, then for every $\varepsilon\in(0,1)$ and $C>1$ there exists $X(\varepsilon,C)\ge 1$ such that, for all $x\ge X(\varepsilon,C)$ and $y>0$,
\begin{equation}\label{eq:potter-EN-eps}
\frac{L(xy)}{L(x)}\ \le\ C\,\max\{y^\varepsilon,y^{-\varepsilon}\}.
\end{equation}
Moreover, for every compact interval $0<a\le y\le b<\infty$,
\begin{equation}\label{eq:UCT-EN-eps}
\lim_{x\to\infty}\ \sup_{y\in[a,b]}\ \left|\frac{L(xy)}{L(x)}-1\right|=0.
\end{equation}
\end{lemma}

\begin{lemma}\label{lem:convex-ineq-EN-eps}
For $x_i\ge 0$ and $\gamma> 1$,
\begin{equation}\label{eq:superadd}
\Big(\sum_{i=1}^n x_i\Big)^{\gamma}\ \ge\ \sum_{i=1}^n x_i^{\gamma}.
\end{equation}
Hence, if $T_i$ are i.i.d.\ nonnegative random variables and $a\ge0$,
\begin{equation}\label{eq:prod-bound-EN-eps}
\EE\!\big[e^{-a(\sum_{i=1}^n T_i)^\gamma}\big]
\ \le\ \EE\!\big[e^{-a\sum_{i=1}^n T_i^\gamma}\big]
\ =\ \prod_{i=1}^n \EE\!\big[e^{-a T_i^\gamma}\big].
\end{equation}
\end{lemma}

\begin{lemma}\label{lem:tauber-Tgamma-EN-eps}
Let $\gamma>1$. Then for every $\varepsilon\in(0,1)$ there exist $a_\varepsilon\in(0,1)$ and positive constants $c_1(\varepsilon)\le c_2(\varepsilon)$ such that, for all $a\in(0,a_\varepsilon]$,
\begin{equation}\label{eq:Lap-two-sided-EN-eps}
c_1(\varepsilon)\,a^{1/(2\gamma)}\ \le\ 1-\EE\!\big[e^{-a\,\tau_1^\gamma}\big]\ \le\ c_2(\varepsilon)\,a^{1/(2\gamma)}.
\end{equation}
Moreover,
\begin{equation}\label{eq:Lap-limit-EN-eps}
\lim_{a\downarrow 0}\ \frac{1-\EE(e^{-a \tau_1^\gamma})}{a^{1/(2\gamma)}}\ =\ \Gamma\!\Big(1-\frac{1}{2\gamma}\Big)\sqrt{\frac{2}{\pi}}.
\end{equation}
\end{lemma}

\begin{proof}[Proof of Proposition \ref{prop:upper-gamma>1-EN-eps}]
    Let \(\varepsilon\in \bigl(0,(2\gamma\beta)\wedge\tfrac{1}{2}\bigr)\). By \eqref{sumindp} and Lemmas~\ref{lem:tauber-Tgamma-EN-eps}–\ref{lem:convex-ineq-EN-eps},
there exists $a_\varepsilon>0$ such that, for all $0<a\le a_\varepsilon$,
\[
\EE(e^{-a\,\tau_n^\gamma})\ \le\ \big(\EE(e^{-a \tau_1^\gamma})\big)^n\ \le\ \big(1-c_1(\varepsilon) a^{1/(2\gamma)}\big)^n
\ \le\ \exp\big(-c_1(\varepsilon)\,n\,a^{1/(2\gamma)}\big).
\]
where \(c_1(\varepsilon)=(1-\varepsilon)\Gamma\!\bigl(1-\tfrac{1}{2\gamma}\bigr)\sqrt{\tfrac{2}{\pi}}\). 

Let \(a(1-h)=-\log p\); for \(h\in \bigl[0,\tfrac{1}{2}\bigr]\) one has \(a(1-h)\in [h,2h]\), hence on \(h\in\bigl(0,h_\varepsilon\wedge\tfrac{1}{2}\bigr]\),
\begin{equation}\label{gac1}
\EE\!\big(e^{-a(1-h)\,\tau_n^\gamma}\big)\le \exp\{-c\,n\, h^{1/(2\gamma)}\} 
\end{equation}
with \(c_{\varepsilon}=c_1(\varepsilon)\,2^{-1/(2\gamma)}\). Moreover, by (\ref{eq:edge-band}) there exists \(h_{\varepsilon}\) such that 
\begin{equation}\label{gac2}
f_\pi(1-h)\ \le\
(1+\varepsilon)\,h^{\beta-1}L(1/h),\qquad 0<h\le h_\varepsilon .
\end{equation}

Let \(h_{2,\varepsilon}:=h_{\varepsilon}\wedge\tfrac{a_{\varepsilon}}{2}\wedge\tfrac{1}{2}\). Then
\begin{equation}
    \begin{split}
        \int_{0}^{1}\EE\!\big(p^{\tau_n^{\gamma}}\big)f_{\pi}(p)\,dp
        &=\int_{0}^{1}\EE\!\big(e^{-a(p)\tau_n^{\gamma}}\big)f_{\pi}(p)\,dp\\
        &=\int_{0}^{1-h_{2,\varepsilon}}\EE\!\big(e^{-a(p)\tau_n^{\gamma}}\big)f_{\pi}(p)\,dp+\int_{1-h_{2,\varepsilon}}^{1}\EE\!\big(e^{-a(p)\tau_n^{\gamma}}\big)f_{\pi}(p)\,dp\\
        &=\int_{h_{2,\varepsilon}}^{1}\EE\!\big(e^{-a(1-h)\tau_n^{\gamma}}\big)f_{\pi}(1-h)\,dh+\int_{0}^{h_{2,\varepsilon}}\EE\!\big(e^{-a(1-h)\tau_n^{\gamma}}\big)f_{\pi}(1-h)\,dh\\
        &:=I_{1,n}+I_{2,n}\,.
    \end{split}
\end{equation}

Furthermore, since \(\tau_{n}\ge n\) a.s., one has
\begin{equation}\label{limI_1}
        I_{1,n}\le \int_{h_{2,\varepsilon}}^{1}\EE\!\big(e^{-a(1-h_{2,\varepsilon})\tau_n^{\gamma}}\big)f_{\pi}(1-h)\,dh
        \le e^{-a(1-h_{2,\varepsilon})n^{\gamma}}\int_{h_{2,\varepsilon}}^{1}f_{\pi}(1-h)\,dh
        \xrightarrow[n\to\infty]{}0.
\end{equation}

We now estimate the order of \(I_{2,n}\). By \eqref{gac1}, \eqref{gac2} and making the change of variable \(y=c_{\varepsilon}\,n\,h^{1/(2\gamma)}\),
\begin{equation}
    \begin{split}
        I_{2,n} &\le (1+\varepsilon) \int_{0}^{h_{2,\varepsilon}} e^{-c_{\varepsilon}\,n \,h^{1/(2\gamma)}}h^{\beta-1}L\!\Big(\tfrac{1}{h}\Big)\,dh\\
        &\le 2\gamma\,(1+\varepsilon) \,(n c_{\varepsilon})^{-2\gamma\beta}\int_{0}^{c_{\varepsilon}n h_{2,\varepsilon}}e^{-y}\,y^{2\gamma\beta-1}\,
        L\!\Big(\big(\tfrac{c_{\varepsilon}n}{y}\big)^{2\gamma}\Big)\,dy\,.
    \end{split}
\end{equation}

By Lemma~\ref{lem:potter-EN-eps}, equation (\ref{eq:potter-EN-eps}), given \(C>1\) there exists \(X(\varepsilon,C)\ge 1\) such that, for all \(x\ge X(\varepsilon,C)\) and \(z>0\),
\[
\frac{L(xz)}{L(x)}\ \le\ C\,\max\{z^\varepsilon,z^{-\varepsilon}\}.
\]
Setting \(x=(c_{\varepsilon}n)^{2\gamma}\) and \(z=\tfrac{1}{y}\) we get
\[
\int_{0}^{c_{\varepsilon}n h_{2,\varepsilon}}e^{-y}\,y^{2\gamma\beta-1}\,
L\!\Big(\big(\tfrac{c_{\varepsilon}n}{y}\big)^{2\gamma}\Big)\,dy
\ \le\ 
C\int_{0}^{c_{\varepsilon}n h_{2,\varepsilon}}e^{-y}\,y^{2\gamma\beta-1}\,
L\big((c_{\varepsilon}n)^{2\gamma}\big)\,\max\{y^{\varepsilon},y^{-\varepsilon}\}\,dy,
\]
and by Lemma~\ref{lem:potter-EN-eps}, equation (\ref{eq:UCT-EN-eps}), there exists \(N(\varepsilon,C)\) such that, if \(n\ge N(\varepsilon,C)\),
\begin{equation}\label{limI_2}
\begin{split}
    I_{2,n}&\le 2\gamma\,(1+\varepsilon)^{2} \,(n c_{\varepsilon})^{-2\gamma\beta}L(n^{2\gamma})\,C
    \int_{0}^{c_{\varepsilon}n h_{2,\varepsilon}}e^{-y}\,y^{2\gamma\beta-1}\max\{y^{\varepsilon},y^{-\varepsilon}\}\,dy\\
    &\le 2\gamma\,(1+\varepsilon)^{2} \,(n c_{\varepsilon})^{-2\gamma\beta}L(n^{2\gamma})\,C
    \int_{0}^{\infty}e^{-y}\,y^{2\gamma\beta-1}\max\{y^{\varepsilon},y^{-\varepsilon}\}\,dy\,.
\end{split}
\end{equation}

Therefore, combining \eqref{limI_1} (\(h_{2,\varepsilon} \le \frac{1}{2}\)) and \eqref{limI_2} yields \eqref{eq:upper-gamma>1-EN-eps}.


\end{proof}

\end{document}
